\begin{table}[H]
\centering
\caption{COCA frequencies for expanded quasi-count nouns}
\label{tab:expansion}
\begin{tabular}{@{}l r r r r r@{}}
\toprule
Noun & COCA freq. & \mention{three} & \mention{several} & \mention{many} & \mention{how many} \\
\midrule
\mention{youth}     & 49{,}846 & 10 & 22 & 86 & 3 \\
\mention{clergy}    & 5{,}879  & 3  & 4  & 26 & 1 \\
\mention{livestock} & 6{,}982  & 1  & 1  & 13 & 0 \\
\mention{poultry}   & 4{,}544  & 0  & 1  & 3  & 0 \\
\mention{vermin}    & 1{,}041  & 0  & 0  & 2  & 0 \\
\bottomrule
\end{tabular}

\smallskip
\footnotesize{Note: \mention{Poultry} \mention{three} count corrected from 1 to 0 (the one raw token is a modifier use: \mention{three poultry houses}). \mention{Clergy} tokens are confirmed head uses. \mention{Livestock} and \mention{youth} counts are unfiltered; \mention{livestock} functions freely as a modifier (26.7\% modifier rate), so its head-use count may be lower. Some \mention{three youth} tokens may reflect the count-singular sense (\mention{a youth} = a young person). Filtering can only reduce counts, not reverse the implicational ordering. Nouns ordered by acceptance of tight properties. Compare Table~\ref{tab:coca} for core nouns.}
\end{table}
